\documentclass[aspectratio=169,11pt]{ctexbeamer}
\usepackage{amsmath,amssymb,mathtools,bm}
\usepackage{hyperref}
\usepackage{booktabs}
\usepackage{array}
\usepackage{graphicx}

\usetheme{Madrid}
\usecolortheme{default}
\setbeamertemplate{navigation symbols}{}

\definecolor{MainBlue}{RGB}{20,65,130}
\definecolor{AccentRed}{RGB}{180,60,45}
\definecolor{SoftBg}{RGB}{245,248,252}

\newcommand{\e}{\mathrm{e}}
\newcommand{\R}{\mathbb{R}}

% Unified Yingxing Beamer footer and visual style.
\newcommand{\huayinglogo}{assets/logo1.jpg}
\newcommand{\huayingslogan}{英行培优，成绩无忧；英行相伴，刮目相看}

\setbeamersize{text margin left=8mm,text margin right=8mm}
\setbeamerfont{frametitle}{size=\large}
\setbeamercolor{structure}{fg=blue!70!black}
\setbeamercolor{frametitle}{bg=blue!75!black,fg=white}
\setbeamercolor{block title}{bg=blue!10,fg=black}
\setbeamercolor{block body}{bg=blue!2,fg=black}
\setbeamercolor{author in head/foot}{bg=blue!8,fg=black}
\setbeamercolor{title in head/foot}{bg=blue!8,fg=black}

\setbeamertemplate{footline}{%
  \leavevmode%
  \hbox{%
    \begin{beamercolorbox}[wd=.12\paperwidth,ht=3.8ex,dp=1.2ex,center]{author in head/foot}%
      \includegraphics[height=2.9ex]{\huayinglogo}%
    \end{beamercolorbox}%
    \begin{beamercolorbox}[wd=.76\paperwidth,ht=3.8ex,dp=1.2ex,center]{title in head/foot}%
      \scriptsize \huayingslogan
    \end{beamercolorbox}%
    \begin{beamercolorbox}[wd=.12\paperwidth,ht=3.8ex,dp=1.2ex,right]{author in head/foot}%
      \scriptsize \insertframenumber/\inserttotalframenumber\hspace*{1em}
    \end{beamercolorbox}%
  }%
}

\title{导数专题讲义}
\subtitle{对数、指数与多极值点最值问题}
\author{}
\date{2026-05-04}

\begin{document}

\begin{frame}
  \titlepage
\end{frame}

\begin{frame}{专题导读}
\small
\begin{block}{核心经验}
含对数时，优先让 $\ln x$ ``单独作战''；含指数时，优先让 $\e^x$ ``寻找搭档''；指数与对数同时出现时，常通过\alert{作商、拆分、变形}来降阶。
\end{block}

\begin{enumerate}
  \item 对数单身狗：把 $\ln x$ 从复杂结构中剥离出来；
  \item 指数找基友：把 $\e^x$ 与代数式组合，转化为新函数；
  \item 指对在一起：常作差、作商或引入辅助函数，避免反复求导。
\end{enumerate}
\end{frame}

\begin{frame}{目录}
  \tableofcontents[sectionstyle=show/shaded,subsectionstyle=show/show/hide]
\end{frame}

\section{专题一：对数与指数型导数问题}
\subsection{对数单身狗}

\begin{frame}{方法总结：对数单身狗}
\small
处理含 $\ln x$ 的不等式或零点问题时，若
\[
F(x)=f(x)\ln x+g(x),
\]
则可将其改写为
\[
\ln x+\frac{g(x)}{f(x)} \quad \text{或} \quad f(x)\left(\ln x+\frac{g(x)}{f(x)}\right).
\]

\begin{block}{优点}
\[
\left(\ln x+\frac{g(x)}{f(x)}\right)'=\frac1x+\left(\frac{g(x)}{f(x)}\right)',
\]
求导后不再含超越函数，单调性与极值分析明显简化。
\end{block}

\begin{alertblock}{适用场景}
参数范围、不等式恒成立、零点个数、极值点判定。
\end{alertblock}
\end{frame}

\begin{frame}{例 1：$(x+1)\ln x-a(x-1)$}
\small
已知
\[
f(x)=(x+1)\ln x-a(x-1).
\]
\begin{enumerate}
  \item 当 $a=4$ 时，求曲线 $y=f(x)$ 在 $(1,f(1))$ 处的切线；
  \item 若 $x\in(1,+\infty)$ 时恒有 $f(x)>0$，求 $a$ 的取值范围。
\end{enumerate}

\pause
\begin{block}{解题主线}
\[
f'(x)=\ln x+\frac{x+1}{x}-a=\ln x+1+\frac1x-a.
\]
令
\[
\phi(x)=\ln x+1+\frac1x-a,
\]
则 $\phi'(x)=\dfrac{x-1}{x^2}$，从而在 $(1,+\infty)$ 上单调递增。
\end{block}
\end{frame}

\begin{frame}{例 1 结论整理}
\small
\textbf{(1) 切线方程}
\[
f(1)=0,\qquad f'(1)=2-a.
\]
当 $a=4$ 时，$f'(1)=-2$，故切线方程为
\[
y=-2(x-1)= -2x+2.
\]

\vspace{0.5em}
\textbf{(2) 参数范围}
\[
f(x)>0 \ (x>1)\Longleftrightarrow f'(x)\ge 0 \text{ 在适当区间控制最小值}.
\]
因为 $\phi(x)$ 在 $(1,+\infty)$ 单调递增，故
\[
\phi(x)>\phi(1)=2-a.
\]
要使 $f(x)$ 在 $(1,+\infty)$ 恒正，只需 $2-a\ge 0$，即
\[
\boxed{a\le 2}.
\]
\end{frame}

\begin{frame}{例 2：切线条件与不等式证明}
\small
设
\[
f(x)=\frac{a\ln x}{x+1}+\frac{b}{x},
\]
且曲线 $y=f(x)$ 在 $(1,f(1))$ 处的切线为
\[
x+2y-3=0.
\]

\begin{block}{第一问}
由切线过点 $(1,1)$ 且斜率为 $-\dfrac12$，得到
\[
f(1)=1,\qquad f'(1)=-\frac12.
\]
联立解得
\[
\boxed{a=1,\ b=1}.
\]
\end{block}
\end{frame}

\begin{frame}{例 2：第二问思路}
\small
代入第一问结果后，
\[
f(x)=\frac{\ln x}{x+1}+\frac1x.
\]

\begin{alertblock}{目标}
证明
\[
f(x)>\frac{\ln x}{x-1}\qquad (x>0,\ x\ne 1)
\]
时，将两边整理成含单个 $\ln x$ 的形式，再研究辅助函数的单调性。
\end{alertblock}
\begin{block}{处理方向}
把不等式移项、通分，转化为
\[
\ln x+\psi(x)>0
\]
的结构，再对新函数求导；这样导数中不再同时出现复杂的对数分式。
\end{block}
\end{frame}

\begin{frame}{对数单身狗：巩固练习}
\small
\begin{enumerate}
  \item 若对任意 $x\ge 1$，
  \[
  x\ln x\ge a(x-1),
  \]
  求实数 $a$ 的取值范围。
  \item 已知
  \[
  f(x)=ax^2-ax-x\ln x,\qquad f(x)\ge 0,
  \]
  求 $a$，并证明 $f(x)$ 存在唯一极大值点。
  \item 已知
  \[
  f(x)=(2+x+ax^2)\ln(1+x)-2x,
  \]
  讨论符号并求极大值点条件下的参数 $a$。
\end{enumerate}
\end{frame}

\subsection{指数找基友}

\begin{frame}{方法总结：指数找基友}
\small
处理含指数函数的不等式时，通常把 $\e^x$ 与代数式\alert{乘除结合}，构造新函数。

\begin{block}{典型变形}
\[
\e^x+f(x)>0
\Longleftrightarrow
1+\frac{f(x)}{\e^x}>0,
\]
\[
\e^x+f(x)=0
\Longleftrightarrow
1+\frac{f(x)}{\e^x}=0.
\]
此时
\[
\left(1+\frac{f(x)}{\e^x}\right)'=\frac{f'(x)-f(x)}{\e^x},
\]
求导后往往退化成多项式或简单有理式。
\end{block}
\end{frame}

\begin{frame}{例 3：$f(x)=\e^x-ax^2$}
\small
\begin{enumerate}
  \item 若 $a=1$，证明：当 $x\ge 0$ 时，$f(x)\ge 1$；
  \item 若 $f(x)$ 在 $(0,+\infty)$ 上只有一个零点，求 $a$。
\end{enumerate}

\pause
\begin{block}{第 1 问}
\[
\e^x-x^2\ge 1
\Longleftrightarrow
(x^2+1)\e^{-x}\le 1.
\]
设
\[
g(x)=(x^2+1)\e^{-x},
\]
则
\[
g'(x)=-(x-1)^2\e^{-x}\le 0.
\]
故 $g(x)\le g(0)=1$，从而
\[
\boxed{\e^x-x^2\ge 1}.
\]
\end{block}
\end{frame}

\begin{frame}{例 3：唯一零点参数}
\small
由
\[
\e^x-ax^2=0 \Longleftrightarrow a=\frac{\e^x}{x^2},
\]
设
\[
\varphi(x)=\frac{\e^x}{x^2},\qquad x>0.
\]
则
\[
\varphi'(x)=\frac{\e^x(x-2)}{x^3}.
\]

\begin{block}{单调性}
\[
\varphi(x)\text{ 在 }(0,2)\text{ 上递减，在 }(2,+\infty)\text{ 上递增}.
\]
\end{block}
\end{frame}

\begin{frame}{例 3：参数结论}
\small
由上一页可知
\[
\varphi_{\min}=\varphi(2)=\frac{\e^2}{4}.
\]

\begin{block}{关键判断}
要使方程在 $(0,+\infty)$ 上只有一个解，必须恰有极小值切到参数线。
\end{block}

因此
\[
\boxed{a=\frac{\e^2}{4}}.
\]
\end{frame}

\begin{frame}{例 4：$f(x)=\e^x+ax^2-x$}
\small
\begin{enumerate}
  \item 当 $a=1$ 时，讨论 $f(x)$ 的单调性；
  \item 当 $x\ge 0$ 时，若
  \[
  f(x)\ge \frac12x^3+1,
  \]
  求 $a$ 的取值范围。
\end{enumerate}

\begin{block}{结论提炼}
\[
f'(x)=\e^x+2x-1.
\]
当 $a=1$ 时，因 $f''(x)=\e^x+2>0$，故 $f'$ 递增，且 $f'(0)=0$，
于是
\[
f(x)\text{ 在 }(-\infty,0)\text{ 递减，在 }(0,+\infty)\text{ 递增}.
\]
\end{block}
\end{frame}

\begin{frame}{例 4：参数范围结果}
\small
将不等式化为
\[
\left(\frac12x^3-ax^2+x+1\right)\e^{-x}\le 1,\qquad x\ge 0.
\]
设
\[
g(x)=\left(\frac12x^3-ax^2+x+1\right)\e^{-x},
\]
可得
\[
g'(x)=-\frac12x(x-2a-1)(x-2)\e^{-x}.
\]

\begin{block}{分区讨论后结论}
通过研究 $2a+1$ 与 $2$ 的位置关系，并结合 $g(0)=1$ 与 $g(2)=(7-4a)\e^{-2}$，得到
\[
\boxed{a\in \left[\frac{7-\e^2}{4},+\infty\right)}.
\]
\end{block}
\end{frame}

\begin{frame}{指数找基友：巩固练习}
\small
\begin{enumerate}
  \item 已知
  \[
  f(x)=\e^x-1-x-ax^2,
  \]
  若 $x\ge 0$ 时恒有 $f(x)\ge 0$，求 $a$ 的取值范围。
  \item 已知
  \[
  f(x)=\e^{-x}+ax\quad (a\in \R),
  \]
  讨论最值；当 $a=0$ 时，证明
  \[
  f(x)>-\frac12x^2+\frac58.
  \]
  \item 已知
  \[
  f(x)=a(x-1),\qquad g(x)=(ax-1)\e^x,
  \]
  讨论相切与不等式整数解问题。
\end{enumerate}
\end{frame}

\subsection{指对分手}

\begin{frame}{方法总结：指对分手}
\small
当 $\e^x$ 与 $\ln x$ 同时出现时，直接求导往往越来越复杂，因此常采用：

\begin{enumerate}
  \item 作商：把 $\e^x\ln x$ 改写成 $\ln x\cdot \e^x$ 的比值或乘积；
  \item 作差：拆成两个单独可控的函数；
  \item 辅助函数：分别研究两边的最值，再比较大小。
\end{enumerate}

\begin{block}{典型想法}
若
\[
F(x)=\e^x\ln x-f(x),
\]
可尝试构造
\[
\frac{F(x)}{\e^x}=\ln x-\frac{f(x)}{\e^x},
\]
这样一边保留 $\ln x$，另一边变成指数找基友的形式。
\end{block}
\end{frame}

\begin{frame}[allowframebreaks]{例 5：$a\e^x\ln x+\dfrac{b\e^{x-1}}{x}$}
\small
已知
\[
f(x)=a\e^x\ln x+\frac{b\e^{x-1}}{x},
\]
其在 $(1,f(1))$ 处的切线为
\[
y=\e(x-1)+2.
\]

\begin{block}{求参}
由切线条件
\[
f(1)=2,\qquad f'(1)=\e,
\]
可解得
\[
\boxed{a=1,\ b=2}.
\]
\end{block}

\pause
\begin{alertblock}{证明不等式的关键变形}
代入后
\[
f(x)=\e^x\ln x+\frac{2\e^{x-1}}{x}.
\]
要证 $f(x)>1$，可化为
\[
x\ln x>x\e^{-x}-\frac2{\e}.
\]
再分别研究
\[
g(x)=x\ln x,\qquad h(x)=x\e^{-x}-\frac2{\e}.
\]
\end{alertblock}
\end{frame}

\begin{frame}{例 5：比较两个辅助函数}
\small
\[
g'(x)=1+\ln x,\qquad h'(x)=\e^{-x}(1-x).
\]

\begin{columns}[t]
\column{0.48\textwidth}
\begin{block}{$g(x)=x\ln x$}
在 $(0,\e^{-1}]$ 上递减，在 $[\e^{-1},+\infty)$ 上递增，
\[
g_{\min}=g(\e^{-1})=-\frac1{\e}.
\]
\end{block}

\column{0.48\textwidth}
\begin{block}{$h(x)=x\e^{-x}-\dfrac2{\e}$}
在 $(0,1]$ 上递增，在 $[1,+\infty)$ 上递减，
\[
h_{\max}=h(1)=-\frac1{\e}.
\]
\end{block}
\end{columns}

\vspace{0.4em}
由于最小值与最大值虽都为 $-\dfrac1{\e}$，但不能在同一点同时取到，故
\[
g(x)>h(x)\quad (x>0),
\]
从而
\[
\boxed{f(x)>1}.
\]
\end{frame}

\begin{frame}{例 6}
\small
\[
f(x)=\frac1x+a\ln x,\qquad g(x)=\frac{\e^x}{x}.
\]
\begin{enumerate}
  \item 讨论函数 $f(x)$ 的单调性；
  \item 当 $a=1$ 时，证明
  \[
  f(x)+g(x)-\left(1+\frac{\e}{x^2}\right)\ln x>\e.
  \]
\end{enumerate}
\end{frame}

\begin{frame}{指对分手：巩固训练}
\small
\begin{enumerate}
  \item 设
  \[
  f(x)=\frac{\ln x+1}{x},
  \]
  证明当 $x>1$ 时，
  \[
  \frac{f(x)}{\e+1}>\frac{2\e^{x-1}}{(x+1)(x\e^x+1)}.
  \]
\end{enumerate}
\end{frame}

\begin{frame}{指对分手：巩固训练续}
\small
\begin{enumerate}
  \item[2.] 已知
  \[
  f(x)=\e^x-a\ln x-a,\quad a>0,
  \]
  当 $a=\e$ 时求极值，并证明
  \[
  \e^{2x-2}-\e^{x-1}\ln x-x\ge 0.
  \]
\end{enumerate}
\end{frame}

\subsection{专题一小结}

\begin{frame}{专题小结}
\small
\begin{center}
\renewcommand{\arraystretch}{1.35}
\begin{tabular}{>{\raggedright\arraybackslash}p{0.22\textwidth} >{\raggedright\arraybackslash}p{0.68\textwidth}}
\toprule
题型特征 & 首选策略 \\
\midrule
含 $\ln x$ 的不等式、零点、参数范围 & 提出或孤立 $\ln x$，构造 $\ln x+\dfrac{g(x)}{f(x)}$ \\
含 $\e^x$ 的恒成立或零点问题 & 让 $\e^x$ 与代数式配对，构造 $1+\dfrac{f(x)}{\e^x}$ \\
同时含 $\e^x,\ln x$ & 作差、作商、分离成两个可控函数比较最值 \\
\bottomrule
\end{tabular}
\end{center}

\begin{alertblock}{最后提醒}
真正决定得分的不是公式记忆，而是看出\alert{该让谁单独作战、该让谁结伴同行、该在什么时候果断分手}。
\end{alertblock}
\end{frame}

\section{专题二：多极值点函数最值问题}

\begin{frame}{专题二导入}
\small
\begin{block}{主题}
处理多极值点函数最值问题的五种途径。
\end{block}

这类问题的共同难点在于：
\begin{enumerate}
  \item 导函数有多个零点，极值点之间存在隐藏关系；
  \item 目标式往往同时含有 $x_1,x_2$ 甚至参数 $a$；
  \item 直接代入计算很乱，必须先\alert{消元}或\alert{降元}。
\end{enumerate}
\end{frame}

\begin{frame}{基本原理}
\small
\begin{enumerate}
  \item 由导函数零点关系消元，构造只含单个极值点 $x_1$ 或 $x_2$ 的函数；
  \item 由韦达关系消元，把目标式转化为参数 $a$ 的函数；
  \item 构造齐次式，引入比值
  \[
  t=\frac{x_2}{x_1}\quad \text{或}\quad t=\frac{x_1}{x_2};
  \]
  \item 使用对数平均不等式；
  \item 利用特殊三变量结构或隐含对称性，例如 $x_1x_2=1$。
\end{enumerate}
\end{frame}

\begin{frame}{五种途径速览}
\small
\begin{columns}[t]
\column{0.48\textwidth}
\begin{block}{途径 1}
构建单个极值点函数：
\[
F(x_1)\ \text{或}\ F(x_2).
\]
\end{block}

\begin{block}{途径 2}
构建参数函数：
\[
G(a).
\]
\end{block}

\column{0.48\textwidth}
\begin{block}{途径 3}
构造齐次式：
\[
t=\frac{x_2}{x_1},\quad t>1.
\]
\end{block}

\begin{block}{途径 4-5}
对数平均不等式；三极值点中的对称与倒数关系。
\end{block}
\end{columns}

\vspace{0.4em}
\begin{alertblock}{判断标准}
看到 $x_1+x_2,\ x_1x_2,\ \ln x_1-\ln x_2,\ \dfrac{x_1}{x_2}$ 这类结构，就该优先考虑消元。
\end{alertblock}
\end{frame}

\begin{frame}{两个高频模型}
\small
原文中特别强调了两个函数模型：
\[
f(x)=x-\frac1x-a\ln x,
\qquad
g(x)=(x+1)\ln x-a(x-1).
\]

\begin{block}{共同特征}
当零点结构合适时，$x=1$ 往往是特殊点；若还存在另外两个零点 $x_1,x_3$，常出现
\[
x_1x_3=1.
\]
\end{block}

\begin{alertblock}{意义}
这类倒数关系可以把双变量问题降成单变量问题，是处理多极值点最值题的核心切口。
\end{alertblock}
\end{frame}

\begin{frame}{例 1：构建单个极值点函数}
\small
已知
\[
f(x)=\frac1x-x+a\ln x.
\]
\begin{enumerate}
  \item 讨论 $f(x)$ 的单调性；
  \item 若 $f(x)$ 存在两个极值点 $x_1,x_2$，证明
  \[
  \frac{f(x_1)-f(x_2)}{x_1-x_2}<a-2.
  \]
\end{enumerate}

\begin{block}{关键关系}
由 $f'(x)=0$ 得
\[
x^2-ax+1=0,
\]
故极值点满足
\[
x_1x_2=1.
\]
\end{block}
\end{frame}

\begin{frame}{例 1：证明思路}
\small
利用 $x_1x_2=1$，原式可化简为只含单个极值点的表达式。

核心是把
\[
\frac{f(x_1)-f(x_2)}{x_1-x_2}
\]
转化为
\[
-2+a\cdot \frac{2\ln x_2}{x_2-\frac1{x_2}}.
\]

进一步只需证明
\[
\frac1{x_2}-x_2+2\ln x_2<0\qquad (x_2>1).
\]

\begin{block}{辅助函数}
设
\[
h(x)=\frac1x-x+2\ln x.
\]
由单调性知 $h(x)$ 在 $(1,+\infty)$ 上递减，且 $h(1)=0$，因此 $h(x_2)<0$。
\end{block}
\end{frame}

\begin{frame}{例 2：构建参数函数}
\small
设
\[
f(x)=ax^2+(2a-1)x+\ln(x+1)
\]
有两个极值点 $x_1,x_2$。

\begin{enumerate}
  \item 求 $a$ 的取值范围；
  \item 证明
  \[
  f(x_1)+f(x_2)<2\ln 2-\frac54.
  \]
\end{enumerate}

\begin{block}{第一问结论}
由 $f'(x)=0$ 转化得到关于 $x+1$ 的二次方程有两个不等正根，从而
\[
\boxed{0<a<\frac18}.
\]
\end{block}
\end{frame}

\begin{frame}[allowframebreaks]{例 2：降到参数}
\small
借助韦达关系，可把
\[
f(x_1)+f(x_2)
\]
完全化为参数函数
\[
f(x_1)+f(x_2)=1-\frac1{4a}-2a-\ln(2a).
\]

令
\[
t=2a,\qquad 0<t<\frac14,
\]
则只需研究
\[
g(t)=1-\frac1{2t}-t-\ln t.
\]

\begin{block}{结论}
由导数分析得 $g(t)$ 单调递增，因此
\[
f(x_1)+f(x_2)<g\!\left(\frac14\right)=2\ln 2-\frac54.
\]
\end{block}
\end{frame}

\begin{frame}[allowframebreaks]{例 3：构造齐次式}
\small
已知
\[
f(x)=2x-a\ln x-\frac1x
\]
有两个不同极值点 $x_1>x_2$。

\begin{enumerate}
  \item 求 $a$ 的取值范围；
  \item 若 $a>3$，求证 $x_1>1$，且
  \[
  \frac{f(x_1)-f(x_2)}{x_1+x_2}<\frac43-2\ln 2.
  \]
\end{enumerate}

\begin{block}{核心变形}
由极值点方程得到
\[
x_1+x_2=\frac a2,\qquad x_1x_2=\frac12.
\]
再令
\[
t=\frac{x_1}{x_2}>2,
\]
目标式被化为单变量函数
\[
h(t)=\frac{4(t-1)}{t+1}-2\ln t.
\]
\end{block}
\end{frame}

\begin{frame}{例 3：单变量收口}
\small
\[
h(t)=\frac{4(t-1)}{t+1}-2\ln t,\qquad t>2.
\]
求导得
\[
h'(t)=\frac{8}{(t+1)^2}-\frac2t<0.
\]

\begin{block}{结论}
$h(t)$ 在 $(2,+\infty)$ 上递减，因此
\[
h(t)<h(2)=\frac43-2\ln 2.
\]
从而原不等式成立。
\end{block}
\end{frame}

\begin{frame}{对数平均不等式法}
\small
若已知极值点满足 $x_1x_2=1$，且目标式中出现
\[
\frac{\ln x_1-\ln x_2}{x_1-x_2},
\]
则可以直接联想到对数平均不等式。

\begin{block}{作用}
它能把含对数差商的表达式直接估住，省掉一轮辅助函数构造。
\end{block}

\begin{alertblock}{记忆点}
当结构里同时出现
\[
x_1x_2=1,\qquad \ln x_1-\ln x_2,\qquad x_1-x_2,
\]
对数平均通常就是捷径。
\end{alertblock}
\end{frame}

\begin{frame}{三极值点与新定义问题}
\small
原文后半部分还讨论了“三极值点问题中的奥秘”和“极值差比系数”等新定义题。

\begin{block}{可提炼的经验}
\begin{enumerate}
  \item 若出现三个极值点，先看导函数是否能拆成三零点结构；
  \item 先找对称中心、倒数关系、参数区间，再决定是否计算；
  \item 新定义问题本质上仍是把极大值与极小值差转化为单参数函数。
\end{enumerate}
\end{block}
\end{frame}

\begin{frame}{专题二小结}
\small
\begin{center}
\renewcommand{\arraystretch}{1.35}
\begin{tabular}{>{\raggedright\arraybackslash}p{0.26\textwidth} >{\raggedright\arraybackslash}p{0.64\textwidth}}
\toprule
结构特征 & 首选方法 \\
\midrule
已知两个极值点满足方程关系 & 韦达消元，转成单个极值点函数 \\
目标式含 $x_1+x_2,\ x_1x_2$ & 参数化，转成 $G(a)$ \\
目标式含比值、对数差、齐次分式 & 设 $t=\dfrac{x_2}{x_1}$ 或 $t=\dfrac{x_1}{x_2}$ \\
出现 $\dfrac{\ln x_1-\ln x_2}{x_1-x_2}$ & 对数平均不等式 \\
出现三个极值点或新定义 & 先画零点结构，再找对称与倒数关系 \\
\bottomrule
\end{tabular}
\end{center}
\end{frame}

\end{document}
